\section{Agronomic Benefits: Yield, Salinity Reclamation, and Socio-Economics}

Beyond its hydrologic and water-quality functions, subsurface drainage is a
production technology whose agronomic and economic value is well established.
The evidence---multi-year paired-field trials, long-term monitoring, and
calibrated simulation across humid, arid, monsoonal, and cold-climate
settings---supports four consolidated conclusions: drainage and its managed
variants raise crop yields, most strongly on poorly drained soils and in wet
years, while maintaining or improving water-use efficiency; subsurface systems
reliably leach salts and reclaim saline and waterlogged land; saline and
brackish water can be reused safely where salt balance is maintained; and
drainage investment delivers favourable financial and social returns. Meta-
analyses anchor the headline magnitudes: a global synthesis of 891 experiments
confirmed drainage significantly reduces soil salinity and raises yield
\parencite{Yang2025}, and a U.S. meta-analysis of 154 comparisons found
controlled drainage and subirrigation raised yield $8.0\%$ over free drainage
\parencite{Singh2022b}. Spacing, depth, and outlet-control mechanisms are
developed in the design section; here we own the yield, reclamation, reuse, and
socio-economic outcomes, together with the salinity-specific design optima and
salt-balance dynamics.

\subsection{Crop yield and agronomic outcomes}

\subsubsection{Magnitude of yield gains}

Direct drained-versus-undrained comparisons repeatedly document substantial
yield increases, primarily by relieving early-season waterlogging. Artificial
drainage raised soybean yield $275$~kg~ha$^{-1}$ ($8\%$) in experiments and
$157$~kg~ha$^{-1}$ ($4\%$) in producer fields \parencite{Mourtzinis2021}, and in
a long-term Indiana experiment corn yields rose progressively in all drained
treatments across 24 corn years while the undrained control stayed flat, though
soybean was unaffected by spacing \parencite{Rui2024,Kladivko2005}. The largest
relative gains occur in heavily stressed environments: Nile-Delta drainage
doubled cotton and rice yields and raised wheat and clover $50\%$
\parencite{ELMOWELHI1988}; Indian acid saline ``kari'' soils gained
$1.36$~t~ha$^{-1}$ ($61\%$) \parencite{Mathew2001}; Bangladeshi waterlogged
saline clay gained $30$--$64\%$ \parencite{Islam2025,Islam2022}; and an Alberta
system raised barley grain yield from $1310$ to $3900$~kg~ha$^{-1}$
\parencite{BENNETT1982}. More moderate but robust gains dominate temperate
systems: shallow drainage in North Carolina raised corn $4\%$ and soybean
$14.3\%$ \parencite{Watkins2024}; Iranian paddies exceeded controls by
$13.4$--$31.3\%$ (up to $66$--$70\%$) \parencite{DarziNaftchali2017b,DarziNaftchali2014};
and Chinese blind-ditch and pipe drainage raised rice $8.0\%$ and $6.0\%$
\parencite{Qi2025b}.

\subsubsection{Controlled drainage, subirrigation, and recycling}

Managed practices deliver benefits beyond free drainage when they retain water
for crop use in dry periods---quantified at $8.0\%$ average yield gain across
154 comparisons \parencite{Singh2022b}. In Missouri, free drainage and drainage
water recycling raised corn yield $20\%$ and $26\%$ over a nondrained control
(recycling reaching $48\%$ in some years), avoiding drought losses
\parencite{Kaur2021}; in North Carolina recycling raised corn and soybean $8\%$
and $17\%$ \parencite{Moursi2024}. Subirrigation is strongly site-dependent,
governed by lateral hydraulic conductivity: across three Ohio sites it improved
relative yield to $80$--$93\%$ at two sites but had almost no effect where
conductivity was very low, and simulations projected it increasingly
outperforming free drainage under future climate \parencite{Gunn2018,Gunn2019}.
Controlled drainage and subirrigation kept soil water optimal in both wet and
dry years \parencite{Niaghi2019}, improved claypan soybean yield while reducing
interannual variability and allowing earlier planting \parencite{Nelson2011},
and combined with foliar pyraclostrobin produced synergistic soybean gains up
to $1070$~kg~ha$^{-1}$ \parencite{Nelson2011b}.

\subsubsection{Mechanisms, water-use efficiency, and boundary conditions}

The proximate mechanism is relief of excess-water stress and restoration of
root-zone aeration \parencite{BENNETT1982,DarziNaftchali2014}, with growth-stage
sensitivity (cotton flowering and boll formation most sensitive; surface
waterlogging worse than subsurface) \parencite{Qian2021}. Improved aeration
propagates through physiology---raising emergence, root vigour, chlorophyll, and
dry matter in cotton \parencite{Yang2022}, lowering leaf Na$^+$/K$^+$ in
sunflower \parencite{Islam2022}, and raising rice nitrogen-use efficiency
$16.1\%$ \parencite{DarziNaftchali2022b}. Controlled drainage repeatedly
improves irrigation water-use efficiency without yield loss: simple management
of existing Egyptian systems improved WUE $15$--$20\%$ \parencite{Wahba2005},
controlled drainage cut applied irrigation $14$--$16\%$ \parencite{Mahmoud2021},
and Chinese buried pipes raised water productivity, highest near $1.4$~m burial
\parencite{Liu2025}. Yield response is soil- and weather-specific: only
undrained controls show negative rainfall responses, wet years favour narrower
spacing and dry years wider \parencite{Rui2024,Kladivko2005,Acharya2019}, and
drainage advanced sowing $7$--$8$ days and avoided $29$--$52\%$ harvest
population losses in wet years \parencite{Mourtzinis2021,Nelson2011}. The
evidence also delineates limits: where surface drainage is adequate, subsurface
drainage did not raise grain yield \parencite{Wiersma2010}; medium-textured
soils are most responsive while fine and coarse show no clear effect
\parencite{Singh2022b}; over-drainage can induce drought
\parencite{Cordeiro2015,DeBoer2006}; an optimum water-table depth exists that
varies with soil, crop, and permeability \parencite{HOOGHOUDT1952,Bigelow2001};
and installation without maintenance underperforms \parencite{Malota2016}.
Adoption is a leading adaptive practice---about three-fourths of surveyed U.S.
producers report drained cropland \parencite{Morton2015}, and nearly $90\%$ of
Finnish farmers rate subsurface drainage important, ahead of irrigation
\parencite{PeltonenSainio2021}.

\subsection{Salt leaching and reclamation of saline and waterlogged soils}

\subsubsection{Reclamation efficacy}

Subsurface drainage reliably lowers the water table below the root zone and
removes salt from the upper profile, converting saline and waterlogged land into
productive cropland, with topsoil salinity reductions commonly $35$--$80\%$ and
corresponding yield gains, validated at production scale
\parencite{Yang2025,Derby2024}. In India, drained-block salinity fell from
$7.18$ to $4.62$~dS~m$^{-1}$ (mean $35.7\%$) while undrained soils rose, and
yields increased $18.75$--$27.62\%$ \parencite{Sharma2006}; wide-spacing
reclamation lowered topsoil salinity from $8.4$ to about $2.1$~dS~m$^{-1}$ and
raised cropping intensity from $143\%$ to $191\%$ \parencite{Manjunatha2004}.
Chinese reclamations are similarly large: drip-plus-pipe cut EC from about
$15$ to $2.9$--$6.9$~dS~m$^{-1}$ and raised seed-cotton yield $52\%$
\parencite{Heng2021}, and after four years the $0$--$1.0$~m layer fell from
$27.6$ to $4.1$--$7.2$~dS~m$^{-1}$ \parencite{Zhang2021}. Desalinization
is surface-concentrated---$0$--$20$~cm rates exceed $80\%$ while deeper layers
lag and the meta-analytic reduction is nonsignificant below $80$~cm
\parencite{Han2023,Yang2025}. Comparable reductions are reported for Turkey,
Tunisia, and Canada, the last raising barley yield up to twentyfold
\parencite{Bahçeci2009,HajAmor2019,Steppuhn2006,Steppuhn2017}, and reclamation
homogenises the field and improves porosity and infiltration
\parencite{Yu2015,Yu2025,Sharma2000}. Where the water table stays shallow,
however, combined waterlogging and salinity are more damaging than waterlogging
alone, with yields failing above $12$~dS~m$^{-1}$ \parencite{Kahlown2002}.

\subsubsection{Design parameters governing salt removal}

Closer spacing and greater depth shorten the desalinization period and increase
salt removal, though depth benefits saturate. A DRAINMOD analysis found salinity
fell $55.5\%$ in two years at $20$~m spacing but only $5.9\%$ in five years at
$200$~m, a slightly-saline target within a decade requiring spacing below $30$~m
\parencite{Chen2025}; the meta-analysis generalises this into climate-specific
layouts (spacing $>30$~m and depth $>1.2$~m in humid climates; $10$--$20$~m
spacing in arid climates with shallow water tables) \parencite{Yang2025}. Field
and column studies identify optima near $15$--$20$~m
\parencite{Wang2020,Han2023,Chen2022}, with desalination decaying beyond about
$30$~m \parencite{Zhu2025b}. Depth effects are pronounced but bounded---pipes
shallower than $1.0$~m are essentially invalid for discharging salt
\parencite{Qian2021b}, yet returns diminish beyond about $1.4$~m
\parencite{Bahçeci2007,Bahçeci2006}, and an intermediate $0.8$--$1.2$~m depth is
often preferred under saline irrigation, saving $10$--$15\%$ water
\parencite{Feng2017,Muhammad2021}. Envelope and filter design is a first-order
control: the envelope factor dominates pipe material, a double-layer geotextile
achieving $>95\%$ salt removal while triple layers clog
\parencite{Li2022,Wang2024d,Elzoghby2021}. Economic analyses converge on roughly
$20$~m spacing at $1.2$~m depth as the maximum-return layout \parencite{Chen2025}.

\subsubsection{Salt balance, temporal decline, and leaching structure}

Because tile drains intercept salt-laden water, effluent salinity is
characteristically higher than irrigation-water salinity and declines over years
as stored salt is depleted \parencite{Li2022,Karegoudar2022}, with ECdw/ECiw
ratios of $11$--$17$ in Iran \parencite{Akram2013}. Salt-balance accounting is
the standard diagnostic: an oasis exported about $94\%$ of imported salt while
discharging $42\%$ of irrigation water \parencite{Askri2022}, though clogging can
cut removal to $15$--$27\%$ \parencite{BenAissa2013,HajAmor2019}. Most residual
salt is removed within ten years, effluent-salinity--drain-age relationships
being stable across decades \parencite{JOHNSTON1993}, though modeling indicates
$15$--$50$ years to full equilibrium, with drainage-water salinity more sensitive
to leaching fraction than to spacing \parencite{Kelleners2000}. Salt export
increases with both depth and spacing, depth having the larger effect where
groundwater quality declines with depth \parencite{GRISMER1993,KAMRA1991,KAMRA1994}.
Leaching is vertically stratified (upper $>$ middle $>$ lower) and ionically
selective, with Cl and Na leached most efficiently ($91\%$ and $88\%$ at
$0$--$20$~cm) and Ca persisting as precipitated gypsum
\parencite{Chen2022,Qian2025,Yu2016}. Re-accumulation is common, following a
``wash-out--recovery--accumulation'' pattern with surface re-salinization in dry
periods \parencite{Zhu2025b,Dou2022,Shi2021b}. Mechanistically, under complete
ponding flow concentrates near the drain ($\sim80\%$ of flow through $20\%$ of
soil volume), so progressive or partial flooding, low-permeability barriers, and
surface mulches raise efficiency substantially (a $10$~m barrier saved
$52$--$82\%$ of leaching water; a mulch cut effluent volume $81\%$)
\parencite{RAO1991,Youngs2000,Yu2018,Zhang2022}, whereas macropore bypass and
fine texture prolong leaching \parencite{Xin2016,MANN1982}.

\subsection{Saline and brackish water reuse}

Subsurface drainage enables safe reuse of saline and brackish water where salt
balance is maintained by adequate leaching. Blended drainage water up to
$9$~dS~m$^{-1}$ produced about $90\%$ wheat yield in India, cyclic use generally
outperforming continuous use \parencite{Sharma2004}; five-year data gave relative
wheat yields of $95.8$, $90.3$, $83.7$, and $77.8\%$ at $6$, $9$, $12$, and
$18.8$~dS~m$^{-1}$ with root-zone salinity kept below $4$~dS~m$^{-1}$
\parencite{Sharma1998}. In China, saline ($6$~g~L$^{-1}$)
irrigation matched freshwater wheat yields with continued salinity decline
\parencite{Gao2023}, and drainage kept saline irrigation viable up to
$6.25$~dS~m$^{-1}$ for maize \parencite{Feng2017} and $8.0$~dS~m$^{-1}$
under cut-soiler drainage \parencite{Neha2022}, with cotton optima at
$2$--$4$~g~L$^{-1}$ \parencite{Li2024,Lu2019}. Reuse nonetheless carries sodicity
and accumulation risks: cyclic reuse caused profile salinity fluctuations and
sodic surface soils (SARe $>11$) \parencite{Mitchell2000,Jayawardane2001},
though modeled reuse over subsurface drains caused only marginal groundwater
deterioration \parencite{KARAJEH1994c}. Integration with complementary practices
raises performance: combined pipe-plus-ditch systems reached $170\%$ of a single
ditch's efficiency \parencite{Yao2024}; timing drainage to snowmelt more than
doubled surface salt leaching \parencite{Yu2016}; cut-soiler shallow drainage cut
salinity $50$--$60\%$ \parencite{Neha2022,Neha2022b}; and regional coupling of
optimised drip irrigation with tailored drainage reduced district-wide irrigation
$10\%$ and salinity $15\%$ \parencite{Liu2026}. At basin scale, persistent
waterlogging and salinity in Central and South Asia reflect deteriorated
infrastructure and inadequate effluent disposal, underscoring that drainage must
be managed within an integrated salt balance
\parencite{Dukhovny2007,Qureshi2008,Ritzema2016}.

\subsection{Socio-economic returns and investment viability}

\subsubsection{Yield, income, and productivity gains}

Drainage investment generates multi-dimensional returns---higher yields, cropping
intensity, farm and household income, land value, and employment
\parencite{Gopalakrishnan2007,AbdelDayem2007}. Drainage typically raises average
yields $10$--$30\%$ and lifts net farm income where leaching water and
sustainable disposal are assured \parencite{AbdelDayem2007,Datta2000,Wichelns2002}.
In the Nile Delta, net incomes rose by $180$--$375$ per hectare with yield gains
accruing progressively from $50\%$ in year one to $100\%$ by year four
\parencite{Wichelns2002}. In India, the Uttar Pradesh project reclaimed over
$175{,}000$~ha of sodic land, raising farm income over $150\%$, cropping
intensity from $49\%$ to over $200\%$, and household income $53\%$
\parencite{Gopalakrishnan2007}, and Haryana pilots raised cropping intensity and
wheat yields as salinity fell from about $50$ to $5$~dS~m$^{-1}$
\parencite{Datta2002,Datta2000}. Long-term drainage also improves the soil
resource base, raising the soil quality index $76$--$136\%$ with carbon-
sequestration co-benefits \parencite{Mukhopadhyay2023,Mukhopadhyay2023b}, and in
Pakistan a project raised cropping intensity from $102\%$ to $157\%$ and living
standards $30\%$ \parencite{Ghumman2011}.

\subsubsection{Financial appraisal, scale, and cost-reducing design}

Formal appraisals confirm favourable performance. The RAJAD project
($14{,}925$~ha) reported a benefit--cost ratio of $1.28$--$2.87$, payback of
$4.5$--$6.9$~years, and an internal rate of return of $18$--$35\%$
\parencite{Gopalakrishnan2007}; Nile-Delta drainage gave farm-level benefit--cost
ratios of $10.3$ (non-saline) and $5.5$ (saline) \parencite{Wichelns2002}, and
small-scale manually installed drainage yielded IRRs of $8$--$11.5\%$ where water
for leaching and disposal are available \parencite{Datta2000,Datta2002}, though
not all appraisals are high \parencite{Ghumman2011}.
National programmes reflect this importance: Egyptian drainage serves about
$2.4$~million hectares \parencite{AbdelDayem2007}, India reclaimed about
$1.39$~million ha \parencite{Gopalakrishnan2007}, and globally $12$--$17\%$ of
cultivated area is drained \parencite{Valipour2014}. Careful design
sharply reduces cost: Nile-Delta design charts gave minimum network costs of
$205$--$374$ per hectare at $300$~m lateral length ($37\%$ below $100$~m)
\parencite{Gabr2024}; corrugated plastic pipe and factory pre-wrapping of
envelopes reduced installation costs by roughly $50\%$ \parencite{Yannopoulos2020};
and mechanization with laser grade control and improved envelopes lowered costs
and raised quality \parencite{Ritzema2006,Ritzema2007,TeKrony2004}. Disposal
infrastructure follows economies of scale, with economic outcomes strongly
conditioned by disposal costs \parencite{Singh2001,Bethune2008,HANSON1994,Napierala2024}.

\subsubsection{Policy and externality context}

Because drainage water carries disposal costs and externalities, its economics
interact with policy. In California an unregulated regime produced applied-water
and drainage volumes above efficient levels, whereas a direct drainage fee
induced efficiency (net social benefit $1132$~US\$/ha)
\parencite{DINAR1989}. Experience underscores that economics alone do not govern
outcomes: management must be technically sound, economically viable, and socially
acceptable \parencite{Letey1999}, and basin-wide accounting shows combination
strategies can cut soil salinity and river salt loading while raising yields
\parencite{Burkhalter2006,Quinn2009,McDowell2021}. Because artificial drainage
shifts diffuse phosphorus toward the subsurface pathway
\parencite{Lucci2014,Allaire2011,Dadfar2010}, and because deeper, wider drains
raise drain-water salinity and salt load, drainage water quality is increasingly
advanced as a necessary design criterion alongside production, with active
water-table management and decision support the key future directions
\parencite{Ayars2015,Yannopoulos2020,Ritzema2023}.
